\noindent\textbf{Аннотация.}
Методом молекулярной динамики исследуется, как предварительная выдержка
сплава Al--Mg--Si с включениями Al$_{13}$Fe$_4$ в постоянном магнитном поле
влияет на пластическую деформацию при одноосной ползучести. Включение
удлинено вдоль поля на 0{,}194\% --- деформация, отвечающая оценённому в
ранних экспериментах межфазному напряжению 147~МПа, --- и удерживается при
этой деформации, пока матрица релаксирует. Разрешённое касательное
напряжение в матрице достигает 15~МПа непосредственно над включением и в
пределах 80~\AA{} спадает до разрешающей способности расчёта. При сдвиговом
нагружении существующая пара дислокаций разрывается начиная с 95--105~МПа, а
новых дислокаций на границе в недеформированной ячейке не образуется вплоть
до 400~МПа. Двухмасштабная оценка воспроизводит наблюдаемый прирост
ползучести 25\% при рассчитанных напряжениях для активационных объёмов,
измеренных для Al--Mg--Si, --- при условии, что включения действительно
деформируются на 0{,}194\%; измеренная магнитострикция сплавов Fe--Al в
двадцать раз меньше, а сохранение эффекта после снятия поля упругим
напряжением не объясняется.

\noindent\textbf{Ключевые слова:} магнитопластичность, магнитная память,
ползучесть, молекулярная динамика, сплавы Al--Mg--Si, Al$_{13}$Fe$_4$,
межфазные напряжения, дислокации, закрепление дислокаций примесями.